\begin{table}[htbp]
\small
\centering
\begin{tabular}{@{}ccccrr@{}}
\toprule
\textbf{Source} & \textbf{Length} & \textbf{Total phase} & \textbf{Cycles} & \textbf{Frequency} & \textbf{At $f_s=512$\,Hz} \\
$i$ & $T_i$ & $5\pi(i+1)$ & $2.5(i+1)$ & [cycles/sample] & [Hz] \\
\midrule
$0$ & $272$ & $5\pi$  & $2.5$ & $0.00923$ & $4.72$ \\
$1$ & $544$ & $10\pi$ & $5.0$ & $0.00921$ & $4.71$ \\
$2$ & $544$ & $15\pi$ & $7.5$ & $0.01381$ & $7.07$ \\
\bottomrule
\end{tabular}
\caption{The sinusoidal pool $\mathcal{P}$ of Light TSMixup, as configured for the probing runs. Source $i$ spans a total phase of $5\pi(i+1)$ radians over its own length $T_i$, so its cycle count is fixed by its index and its frequency by the ratio of the two. The lengths are set from the canonical draw length of $544$ samples, the probing context plus the forecast horizon, which is why two of the three coincide. The last column maps each source onto the sampling rate of the experiments; the pool itself is defined in cycles per sample and is therefore sampling-rate agnostic.}
\label{tab:tsmixupPool}
\end{table}
